\documentclass[11pt]{article}

\usepackage[a4paper,margin=1in]{geometry}
\usepackage{amsmath,amssymb,amsthm,mathtools,bm}
\usepackage{enumitem}
\usepackage{booktabs,tabularx,array}
\usepackage{microtype}
\usepackage[hidelinks]{hyperref}
\usepackage{graphicx}
\usepackage{tikz}
\usetikzlibrary{arrows.meta,positioning,shapes.geometric}

\newtheorem{assumption}{Assumption}[section]
\newtheorem{definition}[assumption]{Definition}
\newtheorem{proposition}[assumption]{Proposition}
\newtheorem{theorem}[assumption]{Theorem}
\newtheorem{lemma}[assumption]{Lemma}
\newtheorem{corollary}[assumption]{Corollary}
\theoremstyle{remark}
\newtheorem{remark}[assumption]{Remark}
\newtheorem{example}[assumption]{Example}

\newcommand{\E}{\mathbb E}
\newcommand{\Pp}{\mathbb P}
\newcommand{\R}{\mathbb R}
\newcommand{\N}{\mathbb N}
\newcommand{\Var}{\operatorname{Var}}
\newcommand{\Cov}{\operatorname{Cov}}
\newcommand{\diag}{\operatorname{diag}}
\newcommand{\op}{\mathrm{o}_{\Pp}}
\newcommand{\Op}{\mathrm{O}_{\Pp}}
\newcommand{\cf}{\mathrm{cf}}
\newcommand{\fact}{\mathrm{fact}}
\newcommand{\cH}{\mathcal H}
\newcommand{\cW}{\mathcal W}
\newcommand{\cT}{\mathcal T}
\newcommand{\cZ}{\mathcal Z}
\newcommand{\cP}{\mathcal P}
\newcommand{\cC}{\mathcal C}
\newcommand{\norm}[1]{\left\lVert #1\right\rVert}
\newcommand{\abs}[1]{\left\lvert #1\right\rvert}
\newcommand{\ip}[2]{\left\langle #1,#2\right\rangle}
\newcommand{\Ph}{\Phi}
\newcommand{\ph}{\phi}
\newcommand{\LamN}{\Lambda_N}
\newcommand{\an}{a_{N,n}}
\newcommand{\bn}{b_{N,n}}
\newcommand{\thetahat}{\widehat\theta_{N,n}}
\newcommand{\sigmahat}{\widehat\sigma_{N,n}}

\title{Work Package 4: Uniform Counterfactual Process Inference for Linear SPDEs\\
\large Mixed-rate functional delta methods, distributional effects, quantiles, tails, and uniform causal hypotheses}
\author{}
\date{7 July 2026}

\begin{document}
\maketitle

\begin{abstract}
This note develops uniform inference for intervention-indexed counterfactual laws generated by a linear diagonal parabolic SPDE.  The target is indexed jointly by an operator intervention, a field-functional weight, and a threshold.  A first correction is conceptual: in the positive Gaussian core the object is a smooth parametric plug-in process in a function space, not a classical empirical process indexed by a Donsker class.  The parameter dimension is two, so all first-order Gaussian limits are finite-rank.  A second correction is statistical: the diffusivity and noise-scale estimators have different rates.  Writing $\an=\sqrt{\sum_{k\le N}\lambda_k}$ and $\bn=\sqrt{Nn}$, with $\an/\bn\to0$, fixed absolute interventions are $\sigma$-only targets and admit a fast $\bn$-rate process, whereas factual laws and factual-versus-counterfactual effects inherit the slower $\an$ rate through the factual diffusivity.  We construct a non-trivial compact infinite-dimensional class of field weights, prove twice continuous Fr\'echet differentiability of the Gaussian CDF map into $C(\cT)$, derive the fast counterfactual process and slow factual/effect processes, and give their covariance kernels.  Quantile and tail processes follow from the same framework.  Supremum statistics for equality and inequality hypotheses are derived, with the contact-set qualification required under composite one-sided nulls.  Finally, we state rate conditions under which Galerkin and temporal approximations are asymptotically negligible.  Multiplier implementation and feasible critical values are deliberately left to Work Package 5.
\end{abstract}

\section{Model, target operators, and mixed-rate input}

Let $H=L^2(D)$ and let $A_0$ be self-adjoint, negative, and have compact resolvent,
\[
 A_0e_k=-\lambda_ke_k,
 \qquad 0<\lambda_1\le\lambda_2\le\cdots,
\]
for an orthonormal basis $(e_k)$.  For $a>0$, write
\[
 S_a(t)=e^{taA_0}.
\]
Let $Q^{1/2}e_k=q_ke_k$, and suppose the stochastic convolution is $H$-valued.  For an absolute operator intervention $do(\theta=a)$, the counterfactual field solves
\begin{equation}\label{eq:cf-spde}
 dU_t^a=aA_0U_t^a\,dt+\sigma Q^{1/2}dW_t,
 \qquad U_0^a=u_0.
\end{equation}
The factual parameter is $(\theta_0,\sigma_0)$, but the intervention overwrites $\theta_0$ by $a$.

Set
\[
 \cH:=L^2(0,T;H),
 \qquad
 \Psi_w(U):=\int_0^T\ip{w_t}{U_t}_H\,dt.
\]
For $w\in\cH$, define the mean and noise-transfer operator
\begin{align}
 m_a(w)
 &:=\int_0^T\ip{w_t}{S_a(t)u_0}_H\,dt,\label{eq:mean}\\
 (K_aw)_k(s)
 &:=q_k\int_s^T e^{-a\lambda_k(t-s)}w_k(t)\,dt,
 \qquad w_k(t):=\ip{w_t}{e_k}_H.\label{eq:K}
\end{align}
Then
\begin{equation}\label{eq:law}
 \Psi_w(U^a)\sim
 \mathcal N\big(m_a(w),\sigma^2V_a(w)\big),
 \qquad
 V_a(w):=\norm{K_aw}_{\cH}^2.
\end{equation}
The fixed-intervention CDF is
\begin{equation}\label{eq:Fcf}
 F^{\cf}_{a,w,z}(\sigma)
 :=\Ph\!\left(c_{a,w,z}(\sigma)\right),
 \qquad
 c_{a,w,z}(\sigma)
 :=\frac{z-m_a(w)}{\sigma\sqrt{V_a(w)}}.
\end{equation}
The factual CDF is the diagonal restriction
\begin{equation}\label{eq:Ffact}
 F^{\fact}_{w,z}(\theta,\sigma)
 :=F^{\cf}_{\theta,w,z}(\sigma).
\end{equation}

Work Package 3 supplies estimators satisfying
\begin{equation}\label{eq:mixed-clt}
 \left(
 \an(\thetahat-\theta_0),
 \bn(\sigmahat-\sigma_0)
 \right)
 \Rightarrow
 (Z_\theta,Z_\sigma),
\end{equation}
where
\begin{equation}\label{eq:rates}
 \an:=\sqrt{\LamN},
 \quad
 \LamN:=\sum_{k=1}^N\lambda_k,
 \quad
 \bn:=\sqrt{Nn},
 \quad
 \frac{\an}{\bn}\to0,
\end{equation}
and
\begin{equation}\label{eq:limit-Z}
 Z_\theta\sim\mathcal N\!\left(0,\frac{2\theta_0}{T}\right),
 \qquad
 Z_\sigma\sim\mathcal N\!\left(0,\frac{\sigma_0^2}{2}\right),
 \qquad Z_\theta\perp Z_\sigma.
\end{equation}

\begin{remark}[Why one scalar rate is incorrect]\label{rem:no-common-rate}
A fixed absolute intervention removes the factual diffusivity from the post-intervention equation.  Hence $F^{\cf}_{a,w,z}$ depends on the estimated factual model only through $\sigma$.  It has the fast $\bn$ rate.  The factual law $F^{\fact}_{w,z}$ depends on $\theta$ and $\sigma$; because $\an/\bn\to0$, its first-order error at the slower $\an$ rate is generated only by $\thetahat$.  No single scalar normalisation gives a non-degenerate joint first-order limit for both target families.
\end{remark}

\section{A non-trivial compact class of field weights}

Let $(\varphi_\ell)_{\ell\ge1}$ be an orthonormal basis of $L^2(0,T)$ and let $\nu_\ell\uparrow\infty$.  Expand
\[
 w(t)=\sum_{k,\ell\ge1}w_{k\ell}\varphi_\ell(t)e_k.
\]
For $r,s>0$ and $R<\infty$, define the ellipsoid
\begin{equation}\label{eq:ellipsoid}
 \cW_R^{r,s}
 :=\left\{
 w\in\cH:
 \sum_{k,\ell\ge1}
 (1+\lambda_k)^{2s}(1+\nu_\ell)^{2r}
 \abs{w_{k\ell}}^2\le R^2
 \right\}.
\end{equation}
Fix compact intervals
\[
 \widetilde\Theta=[\underline a,\overline a]\Subset(0,\infty),
 \qquad
 \cZ=[z_-,z_+],
 \qquad
 \cP=[p_-,p_+]\Subset(0,1).
\]
For constants $0<v_-<v_+<\infty$, set
\begin{equation}\label{eq:W-nondegenerate}
 \cW
 :=\left\{
 w\in\cW_R^{r,s}:
 v_-\le V_a(w)\le v_+
 \text{ for every }a\in\widetilde\Theta
 \right\}.
\end{equation}
The principal index sets are
\begin{align}
 \cT_{\cf}&:=\widetilde\Theta\times\cW\times\cZ,\label{eq:Tcf}\\
 \cT_{\fact}&:=\cW\times\cZ,\label{eq:Tfact}\\
 \cT_Q&:=\widetilde\Theta\times\cW\times\cP.\label{eq:TQ}
\end{align}

\begin{lemma}[Compactness of the weight class]\label{lem:compact-W}
The ellipsoid $\cW_R^{r,s}$ is compact in $\cH$.  If $a\mapsto K_a$ is continuous in operator norm, then $\cW$ is a compact subset of $\cH$ and the index sets in \eqref{eq:Tcf}--\eqref{eq:TQ} are compact metric spaces.
\end{lemma}

\begin{proof}
The embedding of the weighted coefficient space in $\ell^2(\N^2)$ is diagonal.  Since
\[
 (1+\lambda_k)^{2s}(1+\nu_\ell)^{2r}\to\infty
 \quad\text{as }k+\ell\to\infty,
\]
the image of its unit ball has uniformly vanishing coefficient tails and is relatively compact in $\ell^2(\N^2)\simeq\cH$.  It is closed, hence compact.  Continuity of $(a,w)\mapsto V_a(w)=\norm{K_aw}^2$ makes the variance restrictions closed on the compact product $\widetilde\Theta\times\cW_R^{r,s}$.
\end{proof}

\begin{remark}[Finite dictionaries]
A finite scientifically chosen dictionary of space-time weights is an immediate special case.  The ellipsoid is used to show that the theorem is genuinely uniform over an infinite-dimensional target class.  Spatial smoothness alone would not compactify unrestricted temporal oscillations; the temporal weight in \eqref{eq:ellipsoid} is therefore essential.
\end{remark}

\section{Uniform smoothness of the intervention map}

The spectral calculus gives, for $j=1,2$,
\begin{equation}\label{eq:semigroup-derivatives}
 \sup_{a\in\widetilde\Theta}\sup_{t\ge0}
 \norm{t^jA_0^jS_a(t)}_{\mathcal L(H)}<\infty.
\end{equation}
Define
\begin{align}
 \dot m_a(w)
 &:=\int_0^T\ip{w_t}{tA_0S_a(t)u_0}_H\,dt,\label{eq:mdot}\\
 (\dot K_aw)_k(s)
 &:=-q_k\int_s^T\lambda_k(t-s)e^{-a\lambda_k(t-s)}w_k(t)\,dt,\label{eq:Kdot}\\
 \dot V_a(w)
 &:=2\ip{K_aw}{\dot K_aw}_{\cH}.\label{eq:Vdot}
\end{align}
Second derivatives are obtained by replacing the kernels in \eqref{eq:mdot} and \eqref{eq:Kdot} by $t^2A_0^2S_a(t)$ and $q_k\lambda_k^2(t-s)^2e^{-a\lambda_k(t-s)}$, respectively.

\begin{assumption}[Noise-transfer regularity]\label{ass:K}
The modal noise sequence satisfies
\[
 \sup_{k\ge1}\frac{q_k}{\lambda_k}<\infty.
\]
This follows, in particular, from the standing $H$-valued convolution condition $\sum_kq_k^2/\lambda_k<\infty$.
\end{assumption}

\begin{lemma}[Twice continuous differentiability]\label{lem:C2}
Under Assumption \ref{ass:K}, the maps
\[
 a\mapsto m_a\in\cH^*,
 \qquad
 a\mapsto K_a\in\mathcal L(\cH),
 \qquad
 (a,w)\mapsto V_a(w)
\]
are twice continuously differentiable on $\widetilde\Theta$.  Their derivatives through order two are uniformly bounded on $\widetilde\Theta\times\cW$.
\end{lemma}

\begin{proof}
For each mode, Young's inequality for the backward Volterra convolution yields
\[
 \norm{(K_aw)_k}_{L^2(0,T)}
 \le \frac{q_k}{a\lambda_k}\norm{w_k}_{L^2(0,T)}.
\]
The $L^1$ norms of the first and second derivative kernels are bounded by $q_k/(a^2\lambda_k)$ and $2q_k/(a^3\lambda_k)$.  Assumption \ref{ass:K} and $a\ge\underline a$ therefore give uniform operator bounds for $K_a$, $\dot K_a$, and $\ddot K_a$.  Dominated convergence proves operator-norm continuity.  The assertions for $m_a$ follow from \eqref{eq:semigroup-derivatives} and Cauchy--Schwarz.  The formulas for $V_a$ follow by the product rule.
\end{proof}

For $c=c_{a,w,z}(\sigma)$, differentiation of \eqref{eq:Fcf} gives
\begin{equation}\label{eq:dsigma}
 D_\sigma F^{\cf}_{a,w,z}(\sigma)
 =-\frac{c\ph(c)}{\sigma}.
\end{equation}
On the factual diagonal, write
\[
 c^{\fact}_{w,z}(\theta,\sigma)
 :=c_{\theta,w,z}(\sigma).
\]
Then
\begin{align}
 D_\theta F^{\fact}_{w,z}(\theta,\sigma)
 &=\ph(c^{\fact})
 \left[
 -\frac{\dot m_\theta(w)}{\sigma\sqrt{V_\theta(w)}}
 -\frac{c^{\fact}\dot V_\theta(w)}{2V_\theta(w)}
 \right],\label{eq:dtheta}\\
 D_\sigma F^{\fact}_{w,z}(\theta,\sigma)
 &=-\frac{c^{\fact}\ph(c^{\fact})}{\sigma}.\label{eq:dsigmafact}
\end{align}

\begin{theorem}[Uniform Fr\'echet differentiability]\label{thm:Frechet}
Let the parameter rectangle be a compact subset of $(0,\infty)^2$ containing $(\theta_0,\sigma_0)$ in its interior.  Under Lemmas \ref{lem:compact-W} and \ref{lem:C2},
\begin{align*}
 \sigma&\longmapsto F^{\cf}(\sigma)
 :=\big(F^{\cf}_{a,w,z}(\sigma)\big)_{(a,w,z)\in\cT_{\cf}}
 \in C(\cT_{\cf}),\\
 (\theta,\sigma)&\longmapsto F^{\fact}(\theta,\sigma)
 :=\big(F^{\fact}_{w,z}(\theta,\sigma)\big)_{(w,z)\in\cT_{\fact}}
 \in C(\cT_{\fact})
\end{align*}
are twice continuously Fr\'echet differentiable.  Their first derivatives are given by \eqref{eq:dsigma}--\eqref{eq:dsigmafact}.  There is a constant $C<\infty$ such that the uniform Taylor remainders satisfy
\begin{align}
 \norm{F^{\cf}(\sigma+h)-F^{\cf}(\sigma)-hD_\sigma F^{\cf}(\sigma)}_\infty
 &\le Ch^2,\label{eq:Taylor-cf}\\
 \norm{F^{\fact}(\vartheta+h)-F^{\fact}(\vartheta)-DF^{\fact}(\vartheta)[h]}_\infty
 &\le C\norm{h}_{\R^2}^2.\label{eq:Taylor-fact}
\end{align}

\end{theorem}

\begin{proof}
Lemma \ref{lem:C2}, compactness, the lower variance bound $v_-$, and the positive lower bound on $\sigma$ imply uniform bounds on the first two derivatives of the standardised Gaussian argument.  The standard normal CDF has bounded derivatives of every finite order.  The Banach-space chain rule gives continuous Fr\'echet differentiability into the spaces of continuous functions, and Taylor's theorem with the uniform second-derivative bound gives \eqref{eq:Taylor-cf}--\eqref{eq:Taylor-fact}.
\end{proof}

\begin{remark}[No entropy integral is needed]\label{rem:no-entropy}
The stochastic input to Theorem \ref{thm:Frechet} is the two-dimensional estimator $(\thetahat,\sigmahat)$.  The target map sends this finite-dimensional vector smoothly into $C(\cT)$.  Tightness therefore follows by the functional delta method and compact-index continuity.  Donsker entropy calculations would be required for a genuinely nonparametric first-stage process, but not for the present parametric positive core.
\end{remark}

\section{Fast fixed-intervention counterfactual process}

Define
\[
 \widehat F^{\cf}_{N,n}(a,w,z)
 :=F^{\cf}_{a,w,z}(\sigmahat),
 \qquad
 F^{\cf}_0(a,w,z)
 :=F^{\cf}_{a,w,z}(\sigma_0).
\]

\begin{theorem}[Uniform fast counterfactual process]\label{thm:fast-cf}
Under \eqref{eq:mixed-clt}--\eqref{eq:limit-Z} and Theorem \ref{thm:Frechet},
\begin{equation}\label{eq:fast-clt}
 \bn\big(\widehat F^{\cf}_{N,n}-F^{\cf}_0\big)
 \Rightarrow \mathbb G_{\cf}
 \quad\text{in }C(\cT_{\cf}),
\end{equation}
where
\begin{equation}\label{eq:Gcf}
 \mathbb G_{\cf}(a,w,z)
 =d_\sigma^{\cf}(a,w,z)Z_\sigma,
 \qquad
 d_\sigma^{\cf}(a,w,z)
 :=-\frac{c_{a,w,z}(\sigma_0)\ph(c_{a,w,z}(\sigma_0))}{\sigma_0}.
\end{equation}
The limit is a continuous rank-one Gaussian process with covariance kernel
\begin{align}
 \Gamma_{\cf}(\tau,\tau')
 &=\frac{\sigma_0^2}{2}
 d_\sigma^{\cf}(\tau)d_\sigma^{\cf}(\tau')\label{eq:Gamma-cf}\\
 &=\frac12
 c_\tau\ph(c_\tau)c_{\tau'}\ph(c_{\tau'}),\nonumber
\end{align}
where $c_\tau=c_{a,w,z}(\sigma_0)$.
\end{theorem}

\begin{proof}
Apply \eqref{eq:Taylor-cf} with $h=\sigmahat-\sigma_0$.  Since $\bn h=\Op(1)$ and $\bn h^2=\Op(\bn^{-1})=\op(1)$, the uniform remainder is negligible.  The derivative is a continuous element of $C(\cT_{\cf})$, so \eqref{eq:mixed-clt} and the continuous mapping theorem give \eqref{eq:fast-clt}.  The covariance follows from \eqref{eq:limit-Z}.
\end{proof}

\begin{corollary}[Policy-to-policy contrasts]\label{cor:policy-contrast}
Fix $a_0\in\widetilde\Theta$ and define
\[
 C_{a,a_0,w,z}(\sigma)
 :=F^{\cf}_{a,w,z}(\sigma)-F^{\cf}_{a_0,w,z}(\sigma).
\]
Uniformly on $\widetilde\Theta\times\cW\times\cZ$,
\[
 \bn\big(\widehat C_{a,a_0,w,z}-C_{a,a_0,w,z}(\sigma_0)\big)
 \Rightarrow
 \big(d_\sigma^{\cf}(a,w,z)-d_\sigma^{\cf}(a_0,w,z)\big)Z_\sigma.
\]
This contrast retains non-trivial intervention-index variation at the fast rate.
\end{corollary}

\section{Slow factual and factual-versus-counterfactual processes}

Define
\[
 \widehat F^{\fact}_{N,n}(w,z)
 :=F^{\fact}_{w,z}(\thetahat,\sigmahat),
 \qquad
 F^{\fact}_0(w,z)
 :=F^{\fact}_{w,z}(\theta_0,\sigma_0).
\]
Let
\begin{equation}\label{eq:dtheta0}
 d_\theta^{\fact}(w,z)
 :=D_\theta F^{\fact}_{w,z}(\theta_0,\sigma_0).
\end{equation}

\begin{theorem}[Uniform slow factual process]\label{thm:slow-fact}
Under the assumptions of Theorem \ref{thm:fast-cf},
\begin{equation}\label{eq:slow-fact-clt}
 \an\big(\widehat F^{\fact}_{N,n}-F^{\fact}_0\big)
 \Rightarrow \mathbb G_{\fact}
 \quad\text{in }C(\cT_{\fact}),
\end{equation}
where
\begin{equation}\label{eq:Gfact}
 \mathbb G_{\fact}(w,z)
 =d_\theta^{\fact}(w,z)Z_\theta.
\end{equation}
Its covariance kernel is
\begin{equation}\label{eq:Gamma-fact}
 \Gamma_{\fact}(\tau,\tau')
 =\frac{2\theta_0}{T}
 d_\theta^{\fact}(\tau)d_\theta^{\fact}(\tau').
\end{equation}
The scale-estimation contribution vanishes because $\an/\bn\to0$.
\end{theorem}

\begin{proof}
Use \eqref{eq:Taylor-fact}.  The linear term is
\[
 d_\theta^{\fact}\an(\thetahat-\theta_0)
 +D_\sigma F^{\fact}\frac{\an}{\bn}\bn(\sigmahat-\sigma_0).
\]
The second term is $\op(1)$ uniformly.  For the remainder,
\[
 \an\norm{(\thetahat-\theta_0,\sigmahat-\sigma_0)}^2
 =\Op(\an^{-1})+\Op(\an/\bn^2)+\Op(\bn^{-1})=\op(1).
\]
The result follows from \eqref{eq:mixed-clt}.
\end{proof}

For $a\in\widetilde\Theta$, define the factual-versus-counterfactual distributional effect
\begin{equation}\label{eq:Delta}
 \Delta_{a,w,z}
 :=F^{\cf}_{a,w,z}(\sigma_0)-F^{\fact}_{w,z}(\theta_0,\sigma_0),
\end{equation}
and its plug-in estimator
\[
 \widehat\Delta_{a,w,z}
 :=F^{\cf}_{a,w,z}(\sigmahat)
 -F^{\fact}_{w,z}(\thetahat,\sigmahat).
\]

\begin{corollary}[Uniform distributional-effect process]\label{cor:effect}
In $C(\widetilde\Theta\times\cW\times\cZ)$,
\begin{equation}\label{eq:effect-clt}
 \an(\widehat\Delta-\Delta)
 \Rightarrow \mathbb G_\Delta,
 \qquad
 \mathbb G_\Delta(a,w,z)
 =-d_\theta^{\fact}(w,z)Z_\theta.
\end{equation}
For absolute interventions, the first-order slow-rate uncertainty is constant in $a$: the post-intervention term is learned at the faster scale and disappears from the slow limit.
\end{corollary}

\begin{remark}[Relative or parameter-linked policies]\label{rem:linked}
To obtain an intervention-indexed slow limit, one may define a known smooth policy map $a=h(\theta,\eta)$, with policy label $\eta$.  Then the counterfactual derivative contains
\[
 \partial_\theta F^{\cf}_{h(\theta,\eta),w,z}
 =\partial_aF^{\cf}_{a,w,z}\big|_{a=h(\theta,\eta)}
 \partial_\theta h(\theta,\eta).
\]
The relative intervention $h(\theta,\eta)=\eta\theta$ is a canonical example.  This is an extension of the absolute-intervention core, not an assumption hidden inside it.
\end{remark}

\section{Quantile and tail processes}

Let $z_p:=\Ph^{-1}(p)$.  The fixed-intervention Gaussian quantile is
\begin{equation}\label{eq:Qcf}
 Q^{\cf}_{a,w}(p;\sigma)
 =m_a(w)+\sigma\sqrt{V_a(w)}z_p.
\end{equation}
The factual quantile is
\[
 Q^{\fact}_w(p;\theta,\sigma)
 =m_\theta(w)+\sigma\sqrt{V_\theta(w)}z_p.
\]

\begin{theorem}[Uniform quantile processes]\label{thm:quantile}
Uniformly over $\cT_Q$,
\begin{equation}\label{eq:Qcf-clt}
 \bn\big(
 Q^{\cf}_{a,w}(p;\sigmahat)
 -Q^{\cf}_{a,w}(p;\sigma_0)
 \big)
 \Rightarrow
 \mathbb Q_{\cf}(a,w,p)
 :=\sqrt{V_a(w)}z_pZ_\sigma.
\end{equation}
Uniformly over $\cW\times\cP$,
\begin{equation}\label{eq:Qfact-clt}
 \an\big(
 Q^{\fact}_w(p;\thetahat,\sigmahat)
 -Q^{\fact}_w(p;\theta_0,\sigma_0)
 \big)
 \Rightarrow
 \mathbb Q_{\fact}(w,p),
\end{equation}
where
\begin{equation}\label{eq:Qfact-limit}
 \mathbb Q_{\fact}(w,p)
 =\left[
 \dot m_{\theta_0}(w)
 +\sigma_0\frac{\dot V_{\theta_0}(w)}{2\sqrt{V_{\theta_0}(w)}}z_p
 \right]Z_\theta.
\end{equation}
Consequently, the factual-versus-counterfactual quantile-effect process has slow limit $-\mathbb Q_{\fact}$.
\end{theorem}

\begin{proof}
The Gaussian formula \eqref{eq:Qcf} is continuously differentiable uniformly on the compact index sets; $p\in\cP\Subset(0,1)$ makes $z_p$ bounded.  Apply the same fast and slow Taylor arguments as in Theorems \ref{thm:fast-cf} and \ref{thm:slow-fact}.  Formula \eqref{eq:Qfact-limit} is the derivative with respect to $\theta$.
\end{proof}

\begin{remark}[Inverse-map interpretation]
For a general non-Gaussian extension, the inverse-CDF derivative is
\[
 h\longmapsto-\frac{h(Q(p))}{f(Q(p))},
\]
requiring a density uniformly bounded away from zero at the indexed quantiles.  In the present Gaussian core, \eqref{eq:Qcf} gives the quantile process directly and the lower variance bound plus $\cP\Subset(0,1)$ provide the same regularity automatically.
\end{remark}

For a threshold $c\in\cZ$, define the tail probability
\[
 R^{\cf}_{a,w}(c;\sigma)
 :=1-F^{\cf}_{a,w,c}(\sigma).
\]

\begin{corollary}[Uniform tail processes]\label{cor:tail}
Every CDF process theorem above transfers to the corresponding tail process by multiplication by $-1$.  In particular,
\[
 \bn(\widehat R^{\cf}-R^{\cf}_0)
 \Rightarrow-\mathbb G_{\cf},
 \qquad
 \an(\widehat R^{\fact}-R^{\fact}_0)
 \Rightarrow-\mathbb G_{\fact}.
\]
No separate tail empirical-process argument is required for thresholds in a fixed compact set.
\end{corollary}

\section{Uniform causal hypotheses and supremum statistics}

Let $\Delta$ be the distributional effect in \eqref{eq:Delta}.  For equality testing over a compact subset $\cT_0\subseteq\widetilde\Theta\times\cW\times\cZ$, define
\begin{equation}\label{eq:Sequality}
 S^{=}_{N,n}
 :=\an\norm{\widehat\Delta}_{\cT_0,\infty}.
\end{equation}

\begin{corollary}[Equality null]\label{cor:equality}
Under
\[
 H_0^{=}:\Delta(\tau)=0\quad\forall\tau\in\cT_0,
\]
\[
 S^{=}_{N,n}
 \Rightarrow
 \sup_{\tau\in\cT_0}\abs{\mathbb G_\Delta(\tau)}.
\]
Because $\mathbb G_\Delta$ is rank one, the limit equals
\[
 \abs{Z_\theta}
 \sup_{(w,z):\exists a,(a,w,z)\in\cT_0}
 \abs{d_\theta^{\fact}(w,z)}.
\]
\end{corollary}

For first-order stochastic dominance in the direction $F^{\cf}\le F^{\fact}$, define
\begin{equation}\label{eq:Sineq}
 S^{\le}_{N,n}
 :=\an\sup_{\tau\in\cT_0}[\widehat\Delta(\tau)]_+.
\end{equation}
Let the contact set be
\[
 \cC_0:=\{\tau\in\cT_0:\Delta(\tau)=0\}.
\]

\begin{theorem}[One-sided composite null and contact set]\label{thm:contact}
Under
\[
 H_0^{\le}:\Delta(\tau)\le0\quad\forall\tau\in\cT_0,
\]
the following statements hold.
\begin{enumerate}[label=(\roman*)]
 \item If $\sup_{\tau\in\cT_0}\Delta(\tau)<0$, then $S^{\le}_{N,n}\to0$ in probability.
 \item If $\cC_0\neq\varnothing$ and the usual local separation condition holds away from $\cC_0$, then
 \[
 S^{\le}_{N,n}
 \Rightarrow
 \sup_{\tau\in\cC_0}[\mathbb G_\Delta(\tau)]_+.
 \]
\end{enumerate}
The same contact-set logic applies to non-positive tail effects and quantile-effect sign restrictions.
\end{theorem}

\begin{proof}[Proof sketch]
Uniform convergence at rate $\an$ implies that strictly negative regions cannot maximise the positive-part statistic asymptotically.  The supremum map is Hadamard directionally differentiable at the null function, with derivative determined by the set of maximisers/contact points.  Applying this derivative to the limit in Corollary \ref{cor:effect} yields the result.
\end{proof}

\begin{remark}[Interface with Work Package 5]
Theorems \ref{thm:fast-cf}--\ref{thm:contact} identify the target Gaussian laws and the correct studentisation rates.  Work Package 5 must construct feasible critical values, estimate contact sets for one-sided hypotheses, and prove conditional validity.  It should not reintroduce a common rate across the fast and slow target families.
\end{remark}

\section{Galerkin and temporal approximation interface}

Suppose numerical target primitives $m_a^{K,\delta}(w)$ and $V_a^{K,\delta}(w)$ are used, where $K$ is a spatial truncation level and $\delta$ a numerical time step.  Define
\begin{equation}\label{eq:eps-num}
 \varepsilon_{K,\delta}
 :=\sup_{a\in\widetilde\Theta,w\in\cW}
 \left(
 \abs{m_a^{K,\delta}(w)-m_a(w)}
 +\abs{V_a^{K,\delta}(w)-V_a(w)}
 \right).
\end{equation}
Assume the approximate variances remain above $v_-/2$ for large $K$ and small $\delta$.

\begin{proposition}[Transfer of numerical error to target laws]\label{prop:numerical}
There is $C<\infty$ such that
\begin{align}
 \norm{F^{K,\delta}-F}_\infty&\le C\varepsilon_{K,\delta},\label{eq:F-num}\\
 \norm{Q^{K,\delta}-Q}_\infty&\le C\varepsilon_{K,\delta}.\label{eq:Q-num}
\end{align}
Consequently:
\begin{enumerate}[label=(\roman*)]
 \item the fast fixed-intervention limits are unchanged if
 \[
 \bn\varepsilon_{K,\delta}\to0;
 \]
 \item the slow factual and effect limits are unchanged if
 \[
 \an\varepsilon_{K,\delta}\to0.
 \]
\end{enumerate}
Consistent plug-in covariance kernels additionally require uniform convergence of the corresponding first derivatives.
\end{proposition}

\begin{proof}
On the region $V\ge v_-/2$ and $\sigma\ge\underline\sigma$, the Gaussian CDF and quantile formulas are uniformly Lipschitz in $(m,V)$.  This gives \eqref{eq:F-num}--\eqref{eq:Q-num}.  Multiplication by the statistical rate proves the two transfer statements.
\end{proof}

If Work Package 1 supplies
\[
 \varepsilon_{K,\delta}
 \lesssim \lambda_{K+1}^{-\alpha}+\delta^\kappa,
\]
then the sufficient undersmoothing conditions are
\begin{align*}
 \bn\big(\lambda_{K+1}^{-\alpha}+\delta^\kappa\big)&\to0
 &&\text{for fast counterfactual targets},\\
 \an\big(\lambda_{K+1}^{-\alpha}+\delta^\kappa\big)&\to0
 &&\text{for factual and effect targets}.
\end{align*}

\section{Completion statement and strategic implication}

The work package establishes uniform weak convergence over a compact, infinite-dimensional class of space-time field weights and compact intervention and threshold sets.  The processes are tight in spaces of continuous functions, and CDF, quantile, tail, and distributional-effect inference are generated by one smooth target calculus.  The decisive structural result is the rate split:
\[
 \boxed{
 \begin{array}{ll}
 \text{fixed absolute intervention targets:}&\bn=\sqrt{Nn},\\[1mm]
 \text{factual and factual-versus-counterfactual targets:}&\an=\sqrt{\sum_{k\le N}\lambda_k}.
 \end{array}}
\]
The positive core does not generate an infinite-rank empirical-process limit.  Its first-order process is finite-rank because all target uncertainty is propagated from the two-dimensional efficient estimator.  A genuinely infinite-rank limit would require an infinite-dimensional estimated nuisance, such as an unknown noise spectrum, observation operator, initial field, or forcing law.  That extension may strengthen a later paper, but it should not be claimed for the present model.

\begin{thebibliography}{9}

\bibitem{vdV1991}
A.~W. van der Vaart.
\newblock On differentiable functionals.
\newblock \emph{The Annals of Statistics}, 19(1):178--204, 1991.
\newblock doi:10.1214/aos/1176347976.

\bibitem{BeutnerZaehle2016}
E.~Beutner and H.~Z\"ahle.
\newblock Functional delta-method for the bootstrap of quasi-Hadamard differentiable functionals.
\newblock \emph{Electronic Journal of Statistics}, 10(1):1181--1222, 2016.
\newblock doi:10.1214/16-EJS1140.

\bibitem{Carcamo2020}
J.~C\'arcamo, A.~Cuevas, and L.-A.~Rodr\'iguez.
\newblock Directional differentiability for supremum-type functionals: Statistical applications.
\newblock \emph{Bernoulli}, 26(3):2143--2175, 2020.
\newblock doi:10.3150/19-BEJ1188.

\bibitem{ZwingmannHolzmann2020}
T.~Zwingmann and H.~Holzmann.
\newblock Weak convergence of quantile and expectile processes under general assumptions.
\newblock \emph{Bernoulli}, 26(1):323--351, 2020.
\newblock doi:10.3150/19-BEJ1127.

\end{thebibliography}

\end{document}
